% BHU PYQ -- Economics: Financial Market (2025-26 NEP)
\documentclass[11pt,a4paper]{article}
\usepackage[a4paper,top=2.5cm,bottom=2.5cm,left=2.8cm,right=2.8cm,headheight=14pt]{geometry}
\usepackage{amsmath,amssymb}
\usepackage{fontspec}
\setmainfont{Kohinoor Devanagari}
\usepackage{microtype,fancyhdr,enumitem,hyperref,lastpage,parskip}

\hypersetup{colorlinks=true,urlcolor=black,linkcolor=black,pdftitle={ECOMV31: Financial Market -- BHU NEP 2025-26}}

\pagestyle{fancy}\fancyhf{}
\renewcommand{\headrulewidth}{0.4pt}\renewcommand{\footrulewidth}{0.4pt}
\fancyhead[L]{\small \textit{Banaras Hindu University}}
\fancyhead[R]{\small \textit{ECOMV31: Financial Market (2025-26)}}
\fancyfoot[C]{\small Page \thepage\ of \pageref{LastPage}}

\newlist{parts}{enumerate}{2}
\setlist[parts,1]{label=(\alph*),leftmargin=*,itemsep=4pt,topsep=3pt}
\setlist[parts,2]{label=(\roman*),leftmargin=*,itemsep=2pt,topsep=2pt}

\newcommand{\pts}[1]{\hfill \textit{[#1]}}

\begin{document}

\begin{center}
  {\large\bfseries BANARAS HINDU UNIVERSITY}\\[2pt]
  {\normalsize B. A. (Hons.) Semester III Examination, 2025-26 (NEP)}\\[6pt]
  \rule{0.8\linewidth}{0.5pt}\\[6pt]
  {\large\bfseries ECONOMICS}\\[4pt]
  {\large\bfseries Paper Code: ECOMV-31 : Financial Market}\\[4pt]
  {\normalsize Time: 2:30 Hours \hfill Full Marks: 70}\\[2pt]
  {\small REF NO. 2025-26/D-III/033}
\end{center}
\rule{\linewidth}{0.4pt}
\smallskip

\noindent \textit{Instructions: This question paper comprises three Sections A, B and C. All questions are compulsory. Write your Roll No.\ at the top immediately on receipt of this question paper.}

\smallskip
\noindent \textit{इस प्रश्न-पत्र में तीन खण्ड अ, ब तथा स हैं। सभी प्रश्न अनिवार्य हैं।}

\bigskip

\section*{SECTION -- A / खण्ड -- अ}
\noindent \textit{Note: Define the following terms in about 50 words each.}\pts{10 $\times$ 2 = 20}

\noindent \textit{निम्नलिखित पदों को लगभग 50 शब्दों में परिभाषित कीजिए।}

\begin{enumerate}[label=\textbf{\arabic*.},leftmargin=*,itemsep=8pt]
  \item Define the following terms:\pts{10 $\times$ 2 = 20}
  \begin{parts}
    \item Money Market / मुद्रा बाजार
    \item Capital Market / पूँजी बाजार
    \item Financial System / वित्तीय प्रणाली
    \item Primary Market / प्राथमिक बाजार
    \item Initial Public Offer (IPO) / प्रारम्भिक सार्वजनिक निर्गम (IPO)
    \item Underlying Assets / अंतर्निहित परिसंपत्तियाँ
    \item Subjective Risk / विषयनिष्ठ (व्यक्तिपरक) जोखिम
    \item Treasury Bills (T-Bills) / ट्रेजरी बिल (टी-बिल)
    \item Currency Swaps / मुद्रा स्वैप
    \item Commodity Derivatives / कमोडिटी डेरिवेटिव (जिंस व्युत्पन्न)
  \end{parts}
\end{enumerate}

\bigskip

\section*{SECTION -- B / खण्ड -- ब}
\noindent \textit{Note: Answer each question in about 100 words. Each question carries 5 marks.}\pts{4 $\times$ 5 = 20}

\noindent \textit{प्रत्येक प्रश्न का उत्तर लगभग 100 शब्दों में दीजिए। प्रत्येक प्रश्न 5 अंकों का है।}

\begin{enumerate}[label=\textbf{\arabic*.},leftmargin=*,start=2,itemsep=12pt]

  \item What are the common types of Financial Derivatives?\pts{5}
  
  वित्तीय व्युत्पन्नों (Derivatives) के सामान्य प्रकार क्या हैं?

  \begin{center}\textbf{OR / अथवा}\end{center}

  Are derivatives adequately regulated in India?\pts{5}
  
  क्या भारत में डेरिवेटिव बाजार पर्याप्त रूप से विनियमित है?

  \item How are Option contracts different from Futures contracts?\pts{5}
  
  ऑप्शन अनुबंध, फ्यूचर्स अनुबंध से किस प्रकार भिन्न हैं?

  \begin{center}\textbf{OR / अथवा}\end{center}

  Under what market conditions is a Put Option exercised?\pts{5}
  
  किन बाजार परिस्थितियों में पुट ऑप्शन (Put Option) का उपयोग किया जाता है?

  \item Is the exchange of principal an essential condition in Currency Swaps?\pts{5}
  
  क्या मुद्रा स्वैप में मूलधन का आदान-प्रदान एक आवश्यक शर्त है?

  \begin{center}\textbf{OR / अथवा}\end{center}

  ``The basis of an interest rate swap is the spread differential.'' Do you agree? Discuss.\pts{5}
  
  ``ब्याज दर स्वैप का आधार स्प्रेड अंतर (Spread Differential) है।'' क्या आप सहमत हैं? विवेचना कीजिए।

  \item Distinguish between Direct Finance and Indirect Finance.\pts{5}
  
  प्रत्यक्ष वित्त तथा अप्रत्यक्ष वित्त में अन्तर स्पष्ट कीजिए।

  \begin{center}\textbf{OR / अथवा}\end{center}

  Distinguish between Organised and Unorganised Money Markets in India.\pts{5}
  
  भारत में संगठित तथा असंगठित मुद्रा बाजारों में अन्तर स्पष्ट कीजिए।

\end{enumerate}

\bigskip

\section*{SECTION -- C / खण्ड -- स}
\noindent \textit{Note: Answer each question in about 250 words. Each question carries 10 marks.}\pts{3 $\times$ 10 = 30}

\noindent \textit{प्रत्येक प्रश्न का उत्तर लगभग 250 शब्दों में दीजिए। प्रत्येक प्रश्न 10 अंकों का है।}

\begin{enumerate}[label=\textbf{\arabic*.},leftmargin=*,start=6,itemsep=14pt]

  \item What is a Financial System? Describe its main components and explain its significance in an economy.\pts{10}
  
  वित्तीय प्रणाली क्या है? इसके मुख्य घटकों का वर्णन कीजिए तथा अर्थव्यवस्था में इसके महत्त्व की व्याख्या कीजिए।

  \begin{center}\textbf{OR / अथवा}\end{center}

  Do forward contracts always lead to profits? Under what circumstances might a trader feel worse off by entering into a forward contract?\pts{10}
  
  क्या फॉरवर्ड अनुबंध हमेशा लाभ की ओर ले जाते हैं? किन परिस्थितियों में एक व्यापारी फॉरवर्ड अनुबंध करके नुकसान महसूस कर सकता है?

  \item In a Credit Default Swap (CDS), what are the common credit trigger events?\pts{10}
  
  क्रेडिट डिफॉल्ट स्वैप (CDS) में सामान्य क्रेडिट ट्रिगर घटनाएँ क्या हैं?

  \begin{center}\textbf{OR / अथवा}\end{center}

  Explain in detail the objectives and major steps involved in the Risk Management process.\pts{10}
  
  जोखिम प्रबंधन प्रक्रिया के उद्देश्यों तथा प्रमुख चरणों की विस्तार से व्याख्या कीजिए।

\end{enumerate}

\end{document}
